\hypertarget{pyshs-gui-user-manual}{%
\section{PySHS GUI --- User Manual}\label{pyshs-gui-user-manual}}

\textbf{PySHS V3.2.x Graphical User Interface}\\
A desktop interface for the Computational Second Harmonic Scattering
codes \textbf{HRS} and \textbf{SHS}

This program is distributed in the hope that it will be useful, but
WITHOUT ANY WARRANTY; without even the implied warranty of
MERCHANTABILITY or FITNESS FOR A PARTICULAR PURPOSE. See the GNU General
Public License for more details.

This user manual accompanies the PySHS Tkinter GUI. It describes how to
install, launch, and operate the interface, and how each menu and input
parameter maps to the underlying C programs documented in the official
PySHS user manual.

To report errors or suggest improvements related to the GUI, contact the
PySHS authors (e.g.~pierre-marie.gassin@enscm.fr) or the maintainer of
this interface.

\begin{center}\rule{0.5\linewidth}{0.5pt}\end{center}

\hypertarget{contents}{%
\subsection{Contents}\label{contents}}

\begin{enumerate}
\def\labelenumi{\arabic{enumi}.}
\tightlist
\item
  Introduction\\
\item
  Requirements and installation\\
\item
  Launching the interface\\
\item
  Overview of the main window\\
\item
  HRS Calculation tab\\
\item
  SHS Calculation tab\\
\item
  Positions / Orientations generation\\
\item
  Results panel and plots\\
\item
  Input file formats\\
\item
  Practical workflow examples\\
\item
  Important notes and troubleshooting\\
\item
  Relation to the command-line PySHS codes
\end{enumerate}

\begin{center}\rule{0.5\linewidth}{0.5pt}\end{center}

\hypertarget{introduction}{%
\subsection{1. Introduction}\label{introduction}}

\hypertarget{purpose-of-the-gui}{%
\subsubsection{1.1 Purpose of the GUI}\label{purpose-of-the-gui}}

PySHS is a computational package for \textbf{Second Harmonic Scattering
(SHS)} calculations. The core engines are written in \textbf{C}
(\texttt{HRS.c}, \texttt{SHS.c}). They are normally run from the
terminal with interactive prompts and text input files.

The \textbf{PySHS GUI} provides a graphical front-end that:

\begin{itemize}
\tightlist
\item
  selects the calculation module (\textbf{HRS} or \textbf{SHS});
\item
  sets all numerical parameters through forms;
\item
  edits or loads the hyperpolarizability tensor (β) and, for SHS, the
  positions/orientations file;
\item
  compiles the C sources if needed;
\item
  runs the calculation and streams progress messages live;
\item
  displays intensity plots and (for SHS) a 3D dipole-assembly plot;
\item
  allows export of results and figures.
\end{itemize}

The scientific meaning of the outputs (coefficients (a\_V), (b\_V),
(c\_V), (I\_2), (I\_4), depolarization ratio, angle-resolved
intensities, etc.) is the same as in the official PySHS manual. This
document focuses on \textbf{how to use the interface}.

\hypertarget{what-the-gui-can-compute}{%
\subsubsection{1.2 What the GUI can
compute}\label{what-the-gui-can-compute}}

\begin{longtable}[]{@{}
  >{\raggedright\arraybackslash}p{(\columnwidth - 4\tabcolsep) * \real{0.1702}}
  >{\raggedright\arraybackslash}p{(\columnwidth - 4\tabcolsep) * \real{0.4255}}
  >{\raggedright\arraybackslash}p{(\columnwidth - 4\tabcolsep) * \real{0.4043}}@{}}
\toprule\noalign{}
\begin{minipage}[b]{\linewidth}\raggedright
Module
\end{minipage} & \begin{minipage}[b]{\linewidth}\raggedright
Physical situation
\end{minipage} & \begin{minipage}[b]{\linewidth}\raggedright
Underlying binary
\end{minipage} \\
\midrule\noalign{}
\endhead
\bottomrule\noalign{}
\endlastfoot
\textbf{HRS} & Incoherent scattering from uncorrelated molecules in
solution (Hyper-Rayleigh Scattering) & \texttt{HRS} \\
\textbf{SHS} & Coherent scattering from N correlated molecules
(supramolecular aggregate) & \texttt{SHS} \\
\end{longtable}

Calculation types (keywords), shared by both modules:

\begin{longtable}[]{@{}
  >{\raggedright\arraybackslash}p{(\columnwidth - 2\tabcolsep) * \real{0.4091}}
  >{\raggedright\arraybackslash}p{(\columnwidth - 2\tabcolsep) * \real{0.5909}}@{}}
\toprule\noalign{}
\begin{minipage}[b]{\linewidth}\raggedright
Keyword
\end{minipage} & \begin{minipage}[b]{\linewidth}\raggedright
Description
\end{minipage} \\
\midrule\noalign{}
\endhead
\bottomrule\noalign{}
\endlastfoot
\texttt{polarplot\_single} & Polarization-resolved intensity at one
fixed scattering angle Θ \\
\texttt{polarplot\_integrate} & Polarization-resolved intensity
integrated between two scattering angles \\
\texttt{angle\_scattering} & Intensity as a function of scattering angle
for fixed polarization combinations \\
\end{longtable}

\begin{center}\rule{0.5\linewidth}{0.5pt}\end{center}

\hypertarget{requirements-and-installation}{%
\subsection{2. Requirements and
installation}\label{requirements-and-installation}}

\hypertarget{software-requirements}{%
\subsubsection{2.1 Software requirements}\label{software-requirements}}

\begin{itemize}
\tightlist
\item
  \textbf{Python 3.8+} with:

  \begin{itemize}
  \tightlist
  \item
    \texttt{tkinter} (usually included; on Debian/Ubuntu:
    \texttt{sudo\ apt\ install\ python3-tk})
  \item
    \texttt{matplotlib} (\texttt{pip\ install\ matplotlib} or system
    package \texttt{python3-matplotlib})
  \item
    \texttt{numpy}
  \end{itemize}
\item
  A \textbf{C compiler}: \texttt{gcc} or \texttt{clang}

  \begin{itemize}
  \tightlist
  \item
    Linux: \texttt{sudo\ apt\ install\ build-essential}\\
  \item
    macOS: \texttt{xcode-select\ -\/-install}\\
  \item
    Windows: MinGW-w64, MSYS2, or equivalent providing \texttt{gcc}
  \end{itemize}
\item
  Optional but recommended on Linux: \texttt{stdbuf} (GNU coreutils) for
  live line-buffered output from the C programs
\end{itemize}

\hypertarget{directory-layout}{%
\subsubsection{2.2 Directory layout}\label{directory-layout}}

Typical layout of the GUI package:

\begin{Shaded}
\begin{Highlighting}[]
\NormalTok{PySHS\_GUI/}
\NormalTok{├── pyshs\_gui.py              \# Main application}
\NormalTok{├── HRS.c / SHS.c             \# C sources}
\NormalTok{├── HRS / SHS                 \# Executables (created by compilation)}
\NormalTok{├── script\_sphere\_random.py}
\NormalTok{├── script\_spheroid\_random.py}
\NormalTok{├── script\_cylinder\_random.py}
\NormalTok{├── graph\_polar.py}
\NormalTok{├── graph\_angle.py}
\NormalTok{├── graph\_assembly.py}
\NormalTok{├── work/                     \# Temporary job directories}
\NormalTok{└── PySHS\_GUI\_User\_Manual.md  \# This manual}
\end{Highlighting}
\end{Shaded}

\hypertarget{compilation}{%
\subsubsection{2.3 Compilation}\label{compilation}}

The GUI can compile the sources automatically (\textbf{Compile HRS.c} /
\textbf{Compile SHS.c} buttons, or at startup).

Manual compilation (equivalent to the official manual):

\begin{Shaded}
\begin{Highlighting}[]
\FunctionTok{gcc} \AttributeTok{{-}o}\NormalTok{ HRS HRS.c }\AttributeTok{{-}lm}
\FunctionTok{gcc} \AttributeTok{{-}o}\NormalTok{ SHS SHS.c }\AttributeTok{{-}lm}
\end{Highlighting}
\end{Shaded}

On Windows, executables are named \texttt{HRS.exe} and \texttt{SHS.exe}.

\begin{center}\rule{0.5\linewidth}{0.5pt}\end{center}

\hypertarget{launching-the-interface}{%
\subsection{3. Launching the interface}\label{launching-the-interface}}

From a terminal, in the GUI directory:

\begin{Shaded}
\begin{Highlighting}[]
\ExtensionTok{python3}\NormalTok{ pyshs\_gui.py}
\end{Highlighting}
\end{Shaded}

On first launch the interface checks for \texttt{gcc}/\texttt{clang} and
compiles \texttt{HRS.c} and \texttt{SHS.c} if the executables are
missing or older than the sources. Status messages appear in the
\textbf{Calculation results} panel at the bottom.

\begin{center}\rule{0.5\linewidth}{0.5pt}\end{center}

\hypertarget{overview-of-the-main-window}{%
\subsection{4. Overview of the main
window}\label{overview-of-the-main-window}}

The window is split into two vertical regions:

\begin{enumerate}
\def\labelenumi{\arabic{enumi}.}
\tightlist
\item
  \textbf{Top --- calculation tabs}

  \begin{itemize}
  \tightlist
  \item
    Tab \textbf{HRS Calculation}\\
  \item
    Tab \textbf{SHS Calculation}
  \end{itemize}
\item
  \textbf{Bottom --- shared results area}

  \begin{itemize}
  \tightlist
  \item
    Left: text log and numerical output\\
  \item
    Right: plots (intensity and, for SHS, dipole assembly)
  \end{itemize}
\end{enumerate}

A horizontal sash between top and bottom can be dragged to give more
space to the calculation forms or to the results.

\begin{center}\rule{0.5\linewidth}{0.5pt}\end{center}

\hypertarget{hrs-calculation-tab}{%
\subsection{5. HRS Calculation tab}\label{hrs-calculation-tab}}

Use this tab for \textbf{incoherent} Hyper-Rayleigh Scattering
(uncorrelated molecules).

\hypertarget{calculation-type}{%
\subsubsection{5.1 Calculation type}\label{calculation-type}}

Dropdown \textbf{Calculation type}:

\begin{longtable}[]{@{}
  >{\raggedright\arraybackslash}p{(\columnwidth - 2\tabcolsep) * \real{0.4375}}
  >{\raggedright\arraybackslash}p{(\columnwidth - 2\tabcolsep) * \real{0.5625}}@{}}
\toprule\noalign{}
\begin{minipage}[b]{\linewidth}\raggedright
Value
\end{minipage} & \begin{minipage}[b]{\linewidth}\raggedright
Meaning
\end{minipage} \\
\midrule\noalign{}
\endhead
\bottomrule\noalign{}
\endlastfoot
\texttt{polarplot\_single} & Polarization plot at one scattering angle
Θ \\
\texttt{polarplot\_integrate} & Polarization plot with integration from
a start angle to an end angle \\
\texttt{angle\_scattering} & Angular distribution of intensity between
0° and 180° \\
\end{longtable}

The parameter fields below change automatically according to the
selected type.

\hypertarget{parameters}{%
\subsubsection{5.2 Parameters}\label{parameters}}

\begin{longtable}[]{@{}
  >{\raggedright\arraybackslash}p{(\columnwidth - 6\tabcolsep) * \real{0.2500}}
  >{\raggedright\arraybackslash}p{(\columnwidth - 6\tabcolsep) * \real{0.1364}}
  >{\raggedright\arraybackslash}p{(\columnwidth - 6\tabcolsep) * \real{0.3182}}
  >{\raggedright\arraybackslash}p{(\columnwidth - 6\tabcolsep) * \real{0.2955}}@{}}
\toprule\noalign{}
\begin{minipage}[b]{\linewidth}\raggedright
Parameter
\end{minipage} & \begin{minipage}[b]{\linewidth}\raggedright
Unit
\end{minipage} & \begin{minipage}[b]{\linewidth}\raggedright
When visible
\end{minipage} & \begin{minipage}[b]{\linewidth}\raggedright
Description
\end{minipage} \\
\midrule\noalign{}
\endhead
\bottomrule\noalign{}
\endlastfoot
\textbf{Scattering angle Θ} & degrees & \texttt{polarplot\_single} &
Laboratory scattering angle. \textbf{0° = forward / transmission};
\textbf{90° = right-angle collection} (common experimental geometry). \\
\textbf{Start angle} & degrees & \texttt{polarplot\_integrate} & Lower
bound of the scattering-angle integration range. \\
\textbf{End angle} & degrees & \texttt{polarplot\_integrate} & Upper
bound of the integration range. \\
\textbf{Number of points} & integer & \texttt{polarplot\_integrate},
\texttt{angle\_scattering} & Number of samples in the angular range (for
integrate: integration points; for angle\_scattering: points from 0° to
180°). \\
\end{longtable}

\hypertarget{hyperpolarizability-tensor-ux3b2}{%
\subsubsection{5.3 Hyperpolarizability tensor
(β)}\label{hyperpolarizability-tensor-ux3b2}}

Text area for the microscopic hyperpolarizability tensor in
\textbf{atomic units (a.u.)}.

\textbf{Required format} (same as the official PySHS manual):

\begin{Shaded}
\begin{Highlighting}[]
\NormalTok{\# comment line (optional but recommended)}
\NormalTok{xxx 0.0}
\NormalTok{xxy 0.0}
\NormalTok{xxz 0.0}
\NormalTok{xyx 0.0}
\NormalTok{xyy 0.0}
\NormalTok{xyz 0.0}
\NormalTok{xzx 0.0}
\NormalTok{xzy 0.0}
\NormalTok{xzz 0.0}
\NormalTok{yxx 0.0}
\NormalTok{yxy 0.0}
\NormalTok{yxz 0.0}
\NormalTok{yyx 0.0}
\NormalTok{yyy 0.0}
\NormalTok{yyz 0.0}
\NormalTok{yzx 0.0}
\NormalTok{yzy 0.0}
\NormalTok{yzz 0.0}
\NormalTok{zxx 0.0}
\NormalTok{zxy 0.0}
\NormalTok{zxz 0.0}
\NormalTok{zyx 0.0}
\NormalTok{zyy 0.0}
\NormalTok{zyz 0.0}
\NormalTok{zzx 0.0}
\NormalTok{zzy 0.0}
\NormalTok{zzz 1.0}
\end{Highlighting}
\end{Shaded}

\begin{itemize}
\tightlist
\item
  Line 1 is treated as a comment by the C program (skipped).\\
\item
  Each following line: \textbf{label} (e.g.~\texttt{zzz}) then
  \textbf{value}.\\
\item
  All 27 Cartesian components of the rank-3 tensor must be present.
\end{itemize}

Buttons:

\begin{itemize}
\tightlist
\item
  \textbf{Load example} --- fills the area with a simple example (only
  (\beta\_\{zzz\} = 1)).\\
\item
  \textbf{Load file\ldots{}} --- loads a β file from disk.
\end{itemize}

\hypertarget{actions}{%
\subsubsection{5.4 Actions}\label{actions}}

\begin{itemize}
\tightlist
\item
  \textbf{▶ Run HRS calculation} --- writes temporary input files, feeds
  interactive answers to \texttt{HRS}, streams stdout to the results
  panel, then loads plots.\\
\item
  \textbf{Compile HRS.c} --- forces recompilation of the HRS binary.
\end{itemize}

\begin{center}\rule{0.5\linewidth}{0.5pt}\end{center}

\hypertarget{shs-calculation-tab}{%
\subsection{6. SHS Calculation tab}\label{shs-calculation-tab}}

Use this tab for \textbf{coherent} Second Harmonic Scattering from a
finite set of correlated molecules (aggregate).

\hypertarget{calculation-type-1}{%
\subsubsection{6.1 Calculation type}\label{calculation-type-1}}

Same keywords as HRS: \texttt{polarplot\_single},
\texttt{polarplot\_integrate}, \texttt{angle\_scattering}.

\hypertarget{parameters-1}{%
\subsubsection{6.2 Parameters}\label{parameters-1}}

\begin{longtable}[]{@{}
  >{\raggedright\arraybackslash}p{(\columnwidth - 6\tabcolsep) * \real{0.2115}}
  >{\raggedright\arraybackslash}p{(\columnwidth - 6\tabcolsep) * \real{0.1154}}
  >{\raggedright\arraybackslash}p{(\columnwidth - 6\tabcolsep) * \real{0.4231}}
  >{\raggedright\arraybackslash}p{(\columnwidth - 6\tabcolsep) * \real{0.2500}}@{}}
\toprule\noalign{}
\begin{minipage}[b]{\linewidth}\raggedright
Parameter
\end{minipage} & \begin{minipage}[b]{\linewidth}\raggedright
Unit
\end{minipage} & \begin{minipage}[b]{\linewidth}\raggedright
Always / conditional
\end{minipage} & \begin{minipage}[b]{\linewidth}\raggedright
Description
\end{minipage} \\
\midrule\noalign{}
\endhead
\bottomrule\noalign{}
\endlastfoot
\textbf{Number of dipoles} & integer & Always shown & Number of
molecules (N) in the aggregate. \textbf{Must equal the number of lines}
in the positions/orientations file. The GUI automatically overwrites
this field with the line count of the orientation data when you run a
calculation (see §11.1). \\
\textbf{Wavelength} & nm & Always & Laser wavelength (\lambda) in
nanometres (e.g.~800). Used to build wave vectors (k = 2\pi n/\lambda)
and (2k). \\
\textbf{Refractive index -- real part} & dimensionless & Always & Real
part of the medium refractive index (n) (e.g.~1.33 for water). \\
\textbf{Refractive index -- imaginary part} & dimensionless & Always &
Imaginary part of (n) (absorption). Use \texttt{0.0} for a non-absorbing
medium. \\
\textbf{Scattering angle Θ} & degrees & \texttt{polarplot\_single} &
Same meaning as in HRS (0° transmission, 90° right angle). \\
\textbf{Start / End angle} & degrees & \texttt{polarplot\_integrate} &
Integration bounds for the scattering angle. \\
\textbf{Number of points} & integer & \texttt{polarplot\_integrate},
\texttt{angle\_scattering} & Number of angular samples. \\
\end{longtable}

The left column is scrollable; the buttons \textbf{Run SHS calculation}
and \textbf{Compile SHS.c} stay pinned at the bottom of that column.

\hypertarget{hyperpolarizability-ux3b2-sub-tab}{%
\subsubsection{6.3 Hyperpolarizability (β)
sub-tab}\label{hyperpolarizability-ux3b2-sub-tab}}

Same format and buttons as in the HRS tab (see §5.3). The tensor is the
\textbf{molecular} hyperpolarizability in the molecular frame; the C
code rotates it into the laboratory frame for each dipole using the
Euler angles from the orientation file.

\hypertarget{positions-orientations-sub-tab}{%
\subsubsection{6.4 Positions / Orientations
sub-tab}\label{positions-orientations-sub-tab}}

This sub-tab holds the mesoscopic description of the aggregate: for each
molecule (j), orientation and position in the mesoscopic frame.

\textbf{File format} (one molecule per line, six columns):

\begin{Shaded}
\begin{Highlighting}[]
\NormalTok{φ\textquotesingle{}   θ\textquotesingle{}   ψ\textquotesingle{}   x\textquotesingle{}   y\textquotesingle{}   z\textquotesingle{}}
\end{Highlighting}
\end{Shaded}

\begin{longtable}[]{@{}
  >{\raggedright\arraybackslash}p{(\columnwidth - 6\tabcolsep) * \real{0.2581}}
  >{\raggedright\arraybackslash}p{(\columnwidth - 6\tabcolsep) * \real{0.2581}}
  >{\raggedright\arraybackslash}p{(\columnwidth - 6\tabcolsep) * \real{0.1935}}
  >{\raggedright\arraybackslash}p{(\columnwidth - 6\tabcolsep) * \real{0.2903}}@{}}
\toprule\noalign{}
\begin{minipage}[b]{\linewidth}\raggedright
Column
\end{minipage} & \begin{minipage}[b]{\linewidth}\raggedright
Symbol
\end{minipage} & \begin{minipage}[b]{\linewidth}\raggedright
Unit
\end{minipage} & \begin{minipage}[b]{\linewidth}\raggedright
Meaning
\end{minipage} \\
\midrule\noalign{}
\endhead
\bottomrule\noalign{}
\endlastfoot
1 & φ′ & \textbf{radian} & Euler angle (ZYZ convention) \\
2 & θ′ & \textbf{radian} & Euler angle \\
3 & ψ′ & \textbf{radian} & Euler angle \\
4--6 & x′, y′, z′ & \textbf{nm} (same length unit as wavelength) &
Position of the molecule \\
\end{longtable}

The number of non-empty valid lines is the number of dipoles (N).

Buttons:

\begin{itemize}
\tightlist
\item
  \textbf{Load file\ldots{}} --- load an existing orientation file;
  updates \textbf{Number of dipoles} automatically.\\
\item
  \textbf{Save as\ldots{}} --- save the current text area to a file.
\end{itemize}

Below these controls is the \textbf{generation} panel (see §7).

\hypertarget{actions-1}{%
\subsubsection{6.5 Actions}\label{actions-1}}

\begin{itemize}
\tightlist
\item
  \textbf{▶ Run SHS calculation} --- runs \texttt{SHS} with the current
  β, orientation data, and parameters.\\
\item
  \textbf{Compile SHS.c} --- recompiles the SHS binary.
\end{itemize}

\begin{center}\rule{0.5\linewidth}{0.5pt}\end{center}

\hypertarget{positions-orientations-generation}{%
\subsection{7. Positions / Orientations
generation}\label{positions-orientations-generation}}

In the SHS tab, under \textbf{Positions / Orientations}, the panel
\textbf{Generate positions / orientations file} builds orientation files
without leaving the GUI.

\hypertarget{script-selection}{%
\subsubsection{7.1 Script selection}\label{script-selection}}

Dropdown \textbf{Choose a script:}

\begin{longtable}[]{@{}
  >{\raggedright\arraybackslash}p{(\columnwidth - 4\tabcolsep) * \real{0.3429}}
  >{\raggedright\arraybackslash}p{(\columnwidth - 4\tabcolsep) * \real{0.3714}}
  >{\raggedright\arraybackslash}p{(\columnwidth - 4\tabcolsep) * \real{0.2857}}@{}}
\toprule\noalign{}
\begin{minipage}[b]{\linewidth}\raggedright
Menu entry
\end{minipage} & \begin{minipage}[b]{\linewidth}\raggedright
Script file
\end{minipage} & \begin{minipage}[b]{\linewidth}\raggedright
Geometry
\end{minipage} \\
\midrule\noalign{}
\endhead
\bottomrule\noalign{}
\endlastfoot
Sphere (script\_sphere\_random.py) & \texttt{script\_sphere\_random.py}
& Dipoles on a sphere, radial orientation + optional Gaussian
disorder \\
Spheroid (script\_spheroid\_random.py) &
\texttt{script\_spheroid\_random.py} & Dipoles on a spheroid (semi-axes
(a), (c)) + disorder \\
Cylinder (script\_cylinder\_random.py) &
\texttt{script\_cylinder\_random.py} & Dipoles on a cylinder surface
grid + disorder \\
Custom script\ldots{} & User-selected \texttt{.py} file & Any script
that writes the 6-column format \\
\end{longtable}

\hypertarget{parameters-per-script}{%
\subsubsection{7.2 Parameters per script}\label{parameters-per-script}}

\textbf{Sphere}

\begin{longtable}[]{@{}
  >{\raggedright\arraybackslash}p{(\columnwidth - 2\tabcolsep) * \real{0.4375}}
  >{\raggedright\arraybackslash}p{(\columnwidth - 2\tabcolsep) * \real{0.5625}}@{}}
\toprule\noalign{}
\begin{minipage}[b]{\linewidth}\raggedright
Field
\end{minipage} & \begin{minipage}[b]{\linewidth}\raggedright
Meaning
\end{minipage} \\
\midrule\noalign{}
\endhead
\bottomrule\noalign{}
\endlastfoot
Radius (nm) & Sphere radius \\
Number of dipoles & (N) \\
σ disorder (rad) & Standard deviation of angular noise around the radial
direction (0 = pure radial) \\
\end{longtable}

\textbf{Spheroid}

\begin{longtable}[]{@{}ll@{}}
\toprule\noalign{}
Field & Meaning \\
\midrule\noalign{}
\endhead
\bottomrule\noalign{}
\endlastfoot
Radius a (nm) & Equatorial semi-axis \\
Radius c (nm) & Polar semi-axis \\
Number of dipoles & (N) \\
σ disorder (rad) & Angular disorder \\
\end{longtable}

\textbf{Cylinder}

\begin{longtable}[]{@{}
  >{\raggedright\arraybackslash}p{(\columnwidth - 2\tabcolsep) * \real{0.4375}}
  >{\raggedright\arraybackslash}p{(\columnwidth - 2\tabcolsep) * \real{0.5625}}@{}}
\toprule\noalign{}
\begin{minipage}[b]{\linewidth}\raggedright
Field
\end{minipage} & \begin{minipage}[b]{\linewidth}\raggedright
Meaning
\end{minipage} \\
\midrule\noalign{}
\endhead
\bottomrule\noalign{}
\endlastfoot
Radius (nm) & Cylinder radius \\
Length (nm) & Cylinder length \\
Dipoles per radius & Number of angular samples around the
circumference \\
Dipoles along length & Number of samples along the axis \\
σ disorder (rad) & Angular disorder \\
\end{longtable}

Total (N =) (dipoles per radius) × (dipoles along length).

\textbf{Custom script}

\begin{itemize}
\tightlist
\item
  \textbf{Browse\ldots{}} --- select a Python script.\\
\item
  \textbf{stdin answers (one per line)} --- answers to the script's
  interactive \texttt{input()} prompts, in order.\\
\item
  The GUI always passes the output file path as \texttt{sys.argv{[}1{]}}
  (same convention as the official PySHS scripts).
\end{itemize}

\hypertarget{generate-button}{%
\subsubsection{7.3 Generate button}\label{generate-button}}

\textbf{▶ Generate positions / orientations file} runs the selected
script, loads the resulting 6-column data into the text area, sets
\textbf{Number of dipoles} to the line count, and refreshes the
\textbf{Dipole assembly} plot.

\begin{center}\rule{0.5\linewidth}{0.5pt}\end{center}

\hypertarget{results-panel-and-plots}{%
\subsection{8. Results panel and plots}\label{results-panel-and-plots}}

\hypertarget{calculation-results-text}{%
\subsubsection{8.1 Calculation results
(text)}\label{calculation-results-text}}

The left bottom panel shows:

\begin{itemize}
\tightlist
\item
  compilation messages;\\
\item
  live stdout/stderr from the C program (progress percentages, angle
  values, \texttt{av}, \texttt{bv}, \texttt{cv}, \ldots);\\
\item
  the content of the main output file written by HRS/SHS;\\
\item
  warnings (e.g.~mismatch between the dipole field and the orientation
  file).
\end{itemize}

Buttons:

\begin{itemize}
\tightlist
\item
  \textbf{Save\ldots{}} --- save the full log to a text file.\\
\item
  \textbf{Clear} --- clear the panel.\\
\item
  \textbf{⏹ Stop calculation} --- kill the running C process (enabled
  only while a job is running).
\end{itemize}

There is \textbf{no time limit} on calculations. Large aggregates ((N
\sim 10\textsuperscript{3)--(10}4)) may run for minutes to hours.

\hypertarget{plots}{%
\subsubsection{8.2 Plots}\label{plots}}

Two sub-tabs:

\textbf{SHS Intensity}

\begin{itemize}
\tightlist
\item
  For \texttt{polarplot\_single} / \texttt{polarplot\_integrate}: curves
  \textbf{IV(γ)} and \textbf{IH(γ)} versus polarization angle γ
  (degrees).\\
\item
  For \texttt{angle\_scattering}: \textbf{IV(0°)}, \textbf{IV(90°)},
  \textbf{IH(0°)}, \textbf{IH(90°)} versus scattering angle (with
  symmetrization about 0° as in the official plotting scripts).
\end{itemize}

\textbf{Dipole assembly} (mainly SHS)

\begin{itemize}
\tightlist
\item
  3D quiver plot of molecular positions and orientation directions,
  built from the orientation data (same role as
  \texttt{graph\_assembly.py} in the official package).
\end{itemize}

Export:

\begin{itemize}
\tightlist
\item
  \textbf{Export intensity plot\ldots{}} --- PNG or PDF.\\
\item
  \textbf{Export assembly plot\ldots{}} --- PNG or PDF.
\end{itemize}

\begin{center}\rule{0.5\linewidth}{0.5pt}\end{center}

\hypertarget{input-file-formats}{%
\subsection{9. Input file formats}\label{input-file-formats}}

\hypertarget{hyperpolarizability-file-input_beta}{%
\subsubsection{\texorpdfstring{9.1 Hyperpolarizability file
(\texttt{input\_beta})}{9.1 Hyperpolarizability file (input\_beta)}}\label{hyperpolarizability-file-input_beta}}

See §5.3. Values in \textbf{atomic units}. First line = comment; then 27
labelled components.

\hypertarget{orientation-position-file-shs-only}{%
\subsubsection{9.2 Orientation / position file (SHS
only)}\label{orientation-position-file-shs-only}}

See §6.4. Angles in \textbf{radians}, positions in \textbf{nm}
(consistent with wavelength in nm). One molecule per line:

\begin{Shaded}
\begin{Highlighting}[]
\NormalTok{phi  theta  psi  x  y  z}
\end{Highlighting}
\end{Shaded}

Example (two molecules):

\begin{Shaded}
\begin{Highlighting}[]
\NormalTok{  0.000000   1.570796   0.000000   10.000000    0.000000    0.000000}
\NormalTok{  3.141593   1.570796   0.000000  {-}10.000000    0.000000    0.000000}
\end{Highlighting}
\end{Shaded}

\begin{center}\rule{0.5\linewidth}{0.5pt}\end{center}

\hypertarget{practical-workflow-examples}{%
\subsection{10. Practical workflow
examples}\label{practical-workflow-examples}}

\hypertarget{hrs-polarization-plot-at-90}{%
\subsubsection{10.1 HRS polarization plot at
90°}\label{hrs-polarization-plot-at-90}}

\begin{enumerate}
\def\labelenumi{\arabic{enumi}.}
\tightlist
\item
  Open tab \textbf{HRS Calculation}.\\
\item
  Set \textbf{Calculation type} = \texttt{polarplot\_single}.\\
\item
  Set \textbf{Scattering angle Θ} = \texttt{90}.\\
\item
  Load or edit the β tensor.\\
\item
  Click \textbf{▶ Run HRS calculation}.\\
\item
  Inspect coefficients in the log (\texttt{av}, \texttt{bv},
  \texttt{cv}, \texttt{I2v}, \texttt{I4v}, \ldots) and the intensity
  plot.
\end{enumerate}

\hypertarget{shs-for-a-spherical-aggregate}{%
\subsubsection{10.2 SHS for a spherical
aggregate}\label{shs-for-a-spherical-aggregate}}

\begin{enumerate}
\def\labelenumi{\arabic{enumi}.}
\tightlist
\item
  Open tab \textbf{SHS Calculation}.\\
\item
  Set \textbf{Calculation type} = \texttt{polarplot\_single}, Θ =
  \texttt{90}.\\
\item
  Set \textbf{Wavelength} = \texttt{800}, \textbf{n real} =
  \texttt{1.33}, \textbf{n imag} = \texttt{0.0}.\\
\item
  Load the molecular β tensor.\\
\item
  In \textbf{Positions / Orientations}, choose \textbf{Sphere}, set
  radius and (N), click \textbf{Generate}.\\
\item
  Confirm that \textbf{Number of dipoles} matches the generated line
  count.\\
\item
  Click \textbf{▶ Run SHS calculation}.\\
\item
  Review intensity plot and dipole assembly plot; export if needed.
\end{enumerate}

\hypertarget{angle-resolved-shs}{%
\subsubsection{10.3 Angle-resolved SHS}\label{angle-resolved-shs}}

\begin{enumerate}
\def\labelenumi{\arabic{enumi}.}
\tightlist
\item
  Set \textbf{Calculation type} = \texttt{angle\_scattering}.\\
\item
  Set \textbf{Number of points} (e.g.~20 → sampling from −180° to +180°
  in the C code's convention).\\
\item
  Provide β and orientation data.\\
\item
  Run and inspect the angle-resolved curves.
\end{enumerate}

\begin{center}\rule{0.5\linewidth}{0.5pt}\end{center}

\hypertarget{important-notes-and-troubleshooting}{%
\subsection{11. Important notes and
troubleshooting}\label{important-notes-and-troubleshooting}}

\hypertarget{number-of-dipoles-must-match-the-orientation-file}{%
\subsubsection{11.1 Number of dipoles must match the orientation
file}\label{number-of-dipoles-must-match-the-orientation-file}}

The C program asks for (N) and then reads \textbf{exactly (N) lines}
from the orientation file.

In coherent SHS, intensities scale approximately as (N\^{}2). If the GUI
field \textbf{Number of dipoles} differs from the number of lines in the
file, results can be wrong by large factors (e.g.~a factor of 400 if (N)
differs by 20).

\textbf{The GUI now forces (N) to the line count of the orientation data
when a calculation is started}, and logs a warning if the field was
inconsistent.

When comparing GUI results with a manual terminal run, use the
\textbf{same} (N), β file, orientation file, wavelength, refractive
index, angle, and keyword.

\hypertarget{large-n-and-memory}{%
\subsubsection{11.2 Large (N) and memory}\label{large-n-and-memory}}

\texttt{SHS.c} allocates large arrays on the stack sized by (N). Very
large (N) (many thousands) may crash or become extremely slow. Prefer
testing with hundreds of dipoles before scaling up.

\hypertarget{live-progress-not-visible}{%
\subsubsection{11.3 Live progress not
visible}\label{live-progress-not-visible}}

Progress lines require line-buffered stdout. The GUI:

\begin{itemize}
\tightlist
\item
  compiles the C sources with \texttt{setvbuf} for line buffering;\\
\item
  launches processes with \texttt{stdbuf\ -oL\ -eL} when available.
\end{itemize}

Some calculation modes print little until the end;
\texttt{angle\_scattering} is usually the most verbose (one message per
angle).

\hypertarget{compiler-not-found}{%
\subsubsection{11.4 Compiler not found}\label{compiler-not-found}}

Install a C compiler and ensure it is on the \texttt{PATH}. Use
\textbf{Compile HRS.c} / \textbf{Compile SHS.c} and read the log panel.

\hypertarget{stopping-a-long-job}{%
\subsubsection{11.5 Stopping a long job}\label{stopping-a-long-job}}

Use \textbf{⏹ Stop calculation} in the results panel. The underlying C
process is terminated.

\begin{center}\rule{0.5\linewidth}{0.5pt}\end{center}

\hypertarget{relation-to-the-command-line-pyshs-codes}{%
\subsection{12. Relation to the command-line PySHS
codes}\label{relation-to-the-command-line-pyshs-codes}}

The GUI does \textbf{not} reimplement the physics. It:

\begin{enumerate}
\def\labelenumi{\arabic{enumi}.}
\item
  writes temporary β (and orientation) files;\\
\item
  builds the same command line as in the official manual, e.g.~

\begin{Shaded}
\begin{Highlighting}[]
\ExtensionTok{./SHS}\NormalTok{ polarplot\_single input\_beta.txt input\_orient.txt output.txt}
\end{Highlighting}
\end{Shaded}
\item
  supplies the interactive answers (\texttt{scanf}) via stdin:

  \begin{itemize}
  \tightlist
  \item
    SHS: (N), wavelength (nm), (\mathrm{Re}(n)), (\mathrm{Im}(n)), then
    angle-related values depending on the keyword;\\
  \item
    HRS: angle-related values only;\\
  \end{itemize}
\item
  reads \texttt{out\_plot} and the main output file for display.
\end{enumerate}

For the definition of laboratory frames, Euler conventions (ZYZ),
intensity formulas, and coefficient definitions ((a\_V), (b\_V), (c\_V),
(I\_2), (I\_4), \ldots), refer to \textbf{Chapters 2--4 of the official
PySHS V3.2.x user manual}.

\hypertarget{equivalent-terminal-commands}{%
\subsubsection{Equivalent terminal
commands}\label{equivalent-terminal-commands}}

\textbf{HRS}

\begin{Shaded}
\begin{Highlighting}[]
\ExtensionTok{./HRS} \OperatorTok{\textless{}}\NormalTok{Keyword}\OperatorTok{\textgreater{}} \OperatorTok{\textless{}}\NormalTok{input\_beta}\OperatorTok{\textgreater{}} \OperatorTok{\textless{}}\NormalTok{outputfile}\OperatorTok{\textgreater{}}
\end{Highlighting}
\end{Shaded}

\textbf{SHS}

\begin{Shaded}
\begin{Highlighting}[]
\ExtensionTok{./SHS} \OperatorTok{\textless{}}\NormalTok{Keyword}\OperatorTok{\textgreater{}} \OperatorTok{\textless{}}\NormalTok{input\_beta}\OperatorTok{\textgreater{}} \OperatorTok{\textless{}}\NormalTok{input\_orientation}\OperatorTok{\textgreater{}} \OperatorTok{\textless{}}\NormalTok{outputfile}\OperatorTok{\textgreater{}}
\end{Highlighting}
\end{Shaded}

Keywords: \texttt{polarplot\_single} \textbar{}
\texttt{polarplot\_integrate} \textbar{} \texttt{angle\_scattering}.

\begin{center}\rule{0.5\linewidth}{0.5pt}\end{center}

\hypertarget{acknowledgements}{%
\subsection{Acknowledgements}\label{acknowledgements}}

PySHS was developed at Institut Charles Gerhardt Montpellier -- ENSCM.\\
Support: CAMOMILS Research National French Agency, ANR-15-CE21-0002.

Contributors to the scientific code include Dr Pierre-Marie Gassin, Dr
Lotfi Boudjema, and Dr Gaelle Gassin. Discussions with
Prof.~Pierre-François Brevet and the precursor work of Dr Julien
Duboisset (HRS\_computing) are gratefully acknowledged.

Please cite:

\begin{quote}
J. Chem. Inf. Model. 2020, 60, 12, 5912--5917
\end{quote}

when publishing results obtained with PySHS (command line or GUI).

\begin{center}\rule{0.5\linewidth}{0.5pt}\end{center}

\emph{End of the PySHS GUI User Manual}
